\section{Result ledger, claim boundary, and the next gate}
\label{sec:tpc63-claims-next}

TPC--63 completes a missing support-provenance interface and proves that
one inherited single-band certificate cannot deliver a fixed-power
exposure saving.  It does not prove that the underlying arithmetic mass is
bad.  The distinction is the main conclusion of this paper.

\subsection{Result ledger}

\begingroup
\small
\begin{longtable}{>{\raggedright\arraybackslash}p{0.23\textwidth}
                 >{\raggedright\arraybackslash}p{0.52\textwidth}
                 >{\raggedright\arraybackslash}p{0.14\textwidth}}
\toprule
Object & Exact conclusion & Status\\
\midrule
\endfirsthead
\toprule
Object & Exact conclusion & Status\\
\midrule
\endhead
Literal provenance
& The upstream retained support is the coefficient-independent projection
  of a full incidence; preservation of already supported pairs does not
  prove exposure on an independently declared native-high carrier.
& L0 logic; L1 audit\\
\addlinespace
Rectangular completion
& Pre-terminal raw-form, terminal-atomic, and grouped-operator faithfulness
  have different finite incidence tests; spurious crosses are explicit.
& L0 iff tests\\
\addlinespace
Canonical pair relation
& When pair-only terminal completion is possible there is a unique least
  terminal-faithful relation; otherwise the full incidence hypergraph is
  the faithful replacement.
& L0 theorem\\
\addlinespace
Maximal eligible core
& Every declared retained relation selects the same native-high carrier as
  its intersection with the coefficient-independent maximal eligible
  relation; rectangular terminal output still needs the faithfulness test.
& L0 identity\\
\addlinespace
Exposure factorization
& Row exposure is exactly the product of parent-to-terminal alignment and
  retained completion inside the aligned measure; the corresponding
  uniform and distinguished constants are sharp minimax quantities.
& L0 theorem; L1 specialization\\
\addlinespace
Single-band geometry
& The exact sandwich is \(R_0\le r<R_{\max}(y)\), up to one fixed-ratio
  transition window.  The balanced band, and every band certified by the
  inherited freezing inequality, is eventually disjoint from the critical
  high face; retuning to \(R_0\asymp J\) makes the inherited freezing loss
  exceed \(1/400\).
& L1 support no-go\\
\addlinespace
Support-only barrier
& After any independently justified retuning, every fully omitted lower
  dyadic block leaves the inherited
  support-cardinality/Brun--Titchmarsh harmonic certificate with no positive
  fixed-power saving, so it cannot pass the strict mass--degree test.
& L1 certificate no-go\\
\addlinespace
Labeled completion
& Orthogonal direct sum over all nonempty dyadic cofactor blocks preserves
  raw energy exactly and uses \(O(\log X)=X^{o(1)}\) labels.
& L0 theorem\\
\addlinespace
Erasure/grouping
& The fiber square masses give the exact operator norm and nonzero singular
  values; fiber equations give the exact kernel and range-restricted
  injectivity test.
& L0 theorem\\
\addlinespace
Endpoint calculus
& Provenance and logarithmic labeling cost zero; every genuine erasure or
  physical loss is charged once, and the complete total must be strictly
  below \(1/400\).
& L0 ledger; L1 use conditional\\
\bottomrule
\end{longtable}
\endgroup

\subsection{The L0/L1/L2 boundary}

\begin{center}
\small
\begin{tabular}{>{\raggedright\arraybackslash}p{0.12\textwidth}
                >{\raggedright\arraybackslash}p{0.38\textwidth}
                >{\raggedright\arraybackslash}p{0.38\textwidth}}
\toprule
Level & Established here & Not licensed by that level\\
\midrule
L0
& Finite incidence identities, canonical/minimal support constructions,
  exact exposure minimax laws, orthogonal block energy, and grouping
  spectrum/kernel tests.
& No assertion about prime distribution, signed M\"obius cancellation, or
  physical fixed-shift reassembly.\\
\addlinespace
L1
& Specialization to one literal fixed nonzero \(h_0\), its actual
  coefficient-free support, raw atomic norm, cofactor window, and the
  failure of the named support-only certificate.
& No positive fixed-power missing-mass estimate, represented lower bound,
  boundary/tail closure, or number-theoretic target lower bound.\\
\addlinespace
L2
& None.
& In particular, no signed fixed-\(h_0\) parity-kernel estimate, fixed
  \(h_0=2\) prime-pair estimate, prime-pair asymptotic, or twin-prime
  implication.\\
\bottomrule
\end{tabular}
\end{center}

The L1 negative result concerns a proof mechanism, not the truth of the
desired weighted estimate.  Calling it an ``exposure barrier'' means that
support cardinality and the inherited harmonic upper bound are
insufficient on the omitted blocks.  It does not mean those blocks carry a
large proportion of the actual raw weight.

\subsection{What remains unproved}

The following inputs remain open.
\begin{enumerate}[leftmargin=2.3em]
 \item The parent-to-terminal alignment ratio is not shown to be positive,
       much less \(1-X^{-\sigma+o(1)}\), on every active physical row.
 \item The single-band low-cofactor complement is not shown to have a
       fixed-power weighted saving in either uniform or distinguished
       scope.
 \item The additional dyadic blocks in the labeled completion are not
       shown automatically to inherit the literal coefficient identity,
       terminal activation, or physical output formula of the original
       band.
 \item The actual physical grouping map and its upstream range are not
       shown to pass the kernel test or to have an endpoint-quality least
       singular value.
 \item There is no single carrier on which critical occupancy, singleton
       survival, represented transfer, alias control, boundary
       reconnection, and the tail estimate are all proved.
 \item No term in the strict endpoint ledger is newly assigned a favorable
       positive margin by this paper.
\end{enumerate}

In particular, setting \(\cI_X=\cI_X^{\max}\) closes only the relation
completion factor.  It cannot retroactively turn a terminal-inactive
parent atom into an active one, and it cannot identify a labeled direct
sum with an unlabeled physical sum.

\subsection{The next theorem: terminal activation in a fixed scope}

The next useful paper should attack the alignment factor before spending a
difficult cancellation estimate on a possibly empty carrier.  There are
two quantitatively different targets.

\paragraph{Positive activation target.}
For a full-coefficient statement, prove with one fixed
\(\lambda_{\rm act}\ge0\)
\begin{equation}
 \inf_{m:W_H(m)>0}\frac{W_A(m)}{W_H(m)}
 \ge X^{-\lambda_{\rm act}+o(1)}.
 \label{eq:tpc63-next-positive-activation}
\end{equation}
For a declared distinguished coefficient vector \(z^*\), the honest
alternative is
\begin{equation}
 \frac{\sum_mW_A(m)|z_m^*|^2}
      {\sum_mW_H(m)|z_m^*|^2}
 \ge X^{-\lambda_{\rm act,*}+o(1)}.
 \label{eq:tpc63-next-distinguished-activation}
\end{equation}
The second statement must not be promoted to the first.

\paragraph{Near-full alignment target.}
The perturbative missing-mass route needs the stronger estimate
\begin{equation}
 \sum_m\bigl(W_H(m)-W_A(m)\bigr)|z_m|^2
 \le X^{-\sigma_{\rm align}+o(1)}
       \sum_mW_H(m)|z_m|^2,
 \qquad \sigma_{\rm align}>0,
 \label{eq:tpc63-next-near-full-alignment}
\end{equation}
in the same uniform or distinguished scope used downstream.  A merely
positive proportion does not imply
\eqref{eq:tpc63-next-near-full-alignment}.

There are then two honest construction routes.
\begin{enumerate}[label=\textbf{Route \arabic*.},leftmargin=4.3em]
 \item \emph{Preserve the single band.}  Prove a genuinely weighted
       terminal or largest-prime-factor estimate on its low-cofactor
       complement.  Reusing only support cardinality and the old harmonic
       bound cannot work by
       \eqref{eq:tpc63-cardinality-certified-saving}.
 \item \emph{Realize the multiscale completion.}  Construct the literal
       coefficient and terminal maps on every added dyadic block, retain
       the block label, and then compute the exact physical erasure on the
       actual upstream range.  The set union and Parseval identity do not
       supply these arrows.
\end{enumerate}

Either route should begin with an exact zero test.  If a required active
row has \(W_A=0<W_H\), the uniform activation route is closed and that
no-go is itself the correct result.  If activation survives, the following
layer should establish a represented-reference Gram lower bound on the
same carrier before returning to determinant collisions and signed alias
cancellation.

\paragraph{Conclusion.}
The provenance gap has now been converted into two separate, testable
questions: does the literal terminal predicate activate enough of the
native parent mass, and can the completed labeled carrier be mapped back
to the physical output without an active-range kernel?  TPC--63 answers
neither by definition.  It proves the exact finite calculus, shows why the
old single-band cardinality proof cannot answer the first question, and
specifies the next activation theorem without crossing the L1/L2 claim
boundary.
